\section{Robustness of the Benchmark Assumptions}
\label{sec:robustness}
The reviewers questioned three modeling choices inherited from the legacy benchmark: the piecewise-linear power curve, Jensen's wake model, and the $4D$ minimum spacing. This section tests how far the conclusions depend on them. Except for the spacing re-optimization in Section~\ref{sec:spacing}, every analysis below \emph{re-evaluates} the 720 stored final layouts of the follow-up experiment (Section~\ref{sec:additional-runlevel}); no layout is re-optimized, so these checks show how sensitive the reported numbers and the algorithm ordering are to the model, not how each algorithm would perform if it were optimized under the alternative model. The scripts and per-run outputs accompany the revision.

\subsection{Independent Evaluator and Its Validation}
\label{sec:validation}
The benchmark objective of Section~\ref{sec:powermodel}, with the wake geometry of Lemmas~1--2, root-sum-square superposition \eqref{eq5}, the Weibull scale transformation \eqref{eq:weibull-scale}, 24 direction bins and 0.5~m/s speed bins, was re-implemented independently of the code that produced the archived results. Applied to the stored coordinates of all 720 follow-up layouts, it reproduces the recorded objective values with a maximum absolute deviation of $8.7\times10^{-11}$ and the recorded wake losses with a maximum deviation of $2.4\times10^{-4}$ objective units. The evaluator is therefore numerically identical to the implementation behind the follow-up data, which also confirms the objective definition and the conversion to expected power and AEP stated in Section~\ref{sec:powermodel}. For Wind Data Set~II, the ideal objective implied by the follow-up data (14630.76 for two turbines) differs by 0.004\% from the historical tabulated value (14631.37), a minor difference that does not affect any conclusion.

\subsection{Nonlinear Power Curve}
\label{sec:powercurve}
To test the linear power segment in \eqref{eq7}, the stored layouts were re-evaluated with the cubic power-curve model~\cite{Carrillo2013}
\[
f_{\rm c}(s)=P_{\text{rated}}\,\frac{s^3-s_{\text{cut-in}}^3}{s_{\text{rated}}^3-s_{\text{cut-in}}^3},\qquad s_{\text{cut-in}}\le s\le s_{\text{rated}},
\]
with $f_{\rm c}=0$ below cut-in and $f_{\rm c}=P_{\text{rated}}$ above rated speed. The curve uses the same cut-in speed, rated speed and rated power as the benchmark, so it isolates the effect of the shape of the power curve. A second variant adds a 25~m/s cut-out. Expected values were computed by quadrature on a 0.01~m/s grid using the same wake-scaled Weibull distributions.

Table~\ref{tab:powercurve} shows that the linear benchmark curve substantially overstates expected power relative to the cubic curve: the free-stream expected power per turbine falls from 936.4 to 773.8~kW for Data Set~I ($-17.4\%$) and from 487.7 to 310.4~kW for Data Set~II ($-36.4\%$), and the farm-level reductions for the stored layouts are $-17.4\%$ to $-18.3\%$ and $-36.4\%$ to $-37.6\%$, respectively. Adding the cut-out mainly affects the high-wind Data Set~I ($-21.4\%$ to $-21.8\%$). The absolute energy values reported in this paper are therefore specific to the benchmark curve and should not be read as turbine-specific AEP estimates. The \emph{relative} comparison is much less sensitive: within each case the Spearman correlation between the linear and cubic objective values of the 120 stored layouts is at least 0.957, and the ordering of the four algorithms by mean expected power is unchanged in five of the six cases. The exception is Data Set~II at 500~m/4, where LX-SSA and PSO differ by only 0.9 objective units under the linear curve (a difference not distinguished by the tests of Table~\ref{tab:exploratory-tests}) and PSO ranks first under the cubic curve.

\subsection{Gaussian Wake Model}
\label{sec:gaussian}
The stored layouts were also re-evaluated with the Gaussian wake model of Bastankhah and Port\'e-Agel~\cite{Bastankhah2014}. For a downstream distance $x>0$ and radial offset $r_\perp$ from the wake centerline,
\[
\delta_{ij}=\left(1-\sqrt{1-\frac{C_T}{8(\sigma/D)^2}}\right)\exp\!\left(-\frac{r_\perp^2}{2\sigma^2}\right),\quad \frac{\sigma}{D}=k^*\frac{x}{D}+0.2\sqrt{\beta},
\]
with $\beta=\tfrac12(1+\sqrt{1-C_T})/\sqrt{1-C_T}$. The deficit is evaluated at the rotor center and combined by the same root-sum-square superposition and Weibull scaling as the Jensen benchmark. The wake growth rate is $k^*=0.04$, with $k^*=0.022$ as a sensitivity value.

Table~\ref{tab:gaussian} shows that the Gaussian model assigns larger mean wake losses to the same layouts in 22 of the 24 algorithm--case combinations (the exceptions are DE in the 750-m and 1000-m cases of Data Set~II). This is expected, because layouts optimized with the Jensen model can exploit its sharp-edged wake cone and the 24 discrete direction bins by placing turbines just outside the modeled wakes, whereas the Gaussian deficit decays smoothly in the lateral direction. The within-case Spearman correlation between Jensen and Gaussian wake losses ranges from 0.38 (Data Set~I, 500~m/4) to 0.93, so the ranking of individual layouts can change noticeably. At the algorithm level, the algorithm with the lowest mean wake loss is the same under both models in five of six cases for $k^*=0.04$ and in all six cases for $k^*=0.022$; the only change is Data Set~II at 500~m/4, where PSO has the lowest Gaussian wake loss for $k^*=0.04$. The full four-algorithm ordering is preserved in four cases for $k^*=0.04$ and in five for $k^*=0.022$. Wake-loss differences between algorithms should therefore be read as properties of the Jensen benchmark; whether they persist when layouts are \emph{optimized} with a higher-fidelity wake model remains to be tested.

\subsection{Minimum-Spacing Sensitivity}
\label{sec:spacing}
The spacing assumption was examined in three complementary ways.

\emph{Geometric capacity.} Table~\ref{tab:capacity} reports, for each radius and spacing, the largest turbine count for which a multi-start packing search constructed a layout satisfying both the circular boundary and the minimum distance. These are constructive lower bounds on the true geometric capacity; the 500-m values (13, 8 and 7) coincide with the known optimal circle-packing results. %%CAPACITY_SENTENCE%%

\emph{Stored layouts.} The layouts optimized under $4D$ are frequently active at the spacing constraint. Of the 240 stored four-turbine layouts, 77.5\% also satisfy $5D$ and 40.4\% satisfy $6D$; of the 480 stored eight-turbine layouts, only 27.3\% satisfy $5D$ and 5.6\% satisfy $6D$ (none of the 750-m, eight-turbine layouts satisfies $6D$). The reported optimal layouts are therefore specific to the $4D$ rule and cannot be transferred to a site with larger spacing requirements without re-optimization.

\emph{Re-optimization.} Because the numerical values of the LX-SSA parameters $\phi$ and $\chi$ are not available in the archive, LX-SSA itself could not be re-run under different spacings without guessing them. The spacing effect was therefore measured with an independent reference implementation of SSA that follows \eqref{eq12} and \eqref{eq14} exactly, uses the boundary projection and feasibility-first rules of Section~\ref{sec:constraints}, a population of 30, 100 iterations (3,030 evaluations, matching the follow-up SSA budget) and 30 independently seeded runs per cell. Table~\ref{tab:spacing-reopt} shows that increasing the spacing from $4D$ to $6D$ has no material effect on the four-turbine 500-m cases (changes smaller than one standard deviation) and a small effect on the 1000-m eight-turbine cases, but it lowers the mean objective of the 750-m eight-turbine cases by 750 units (0.67\%) for Data Set~I and 327 units (0.57\%) for Data Set~II, while the mean wake loss rises from 0.53\% to 1.19\% and from 2.17\% to 2.73\% of the ideal objective. Tighter spacing thus matters most when the farm is densely populated, which is exactly where the historical tables report their largest turbine counts.

\emph{Calibration and budget effect.} The reference SSA also serves as a calibration check. At $4D$ and 3,030 evaluations it attains higher mean objectives than the recorded SSA runs \emph{and} the recorded LX-SSA runs in all six cases (Mann--Whitney $p\le0.008$ in every comparison; e.g., $+339$ and $+543$ units above the recorded LX-SSA mean for the 750-m cases of Data Sets~I and~II). A variant that retains a candidate only if it improves on its incumbent gave essentially the same results. Implementation details that are not documented in the archive (for example initialization, bound handling and the precise leader step) therefore affect the absolute results by more than the recorded LX-SSA--SSA differences. For this reason the reference runs are used only for comparisons \emph{within} one implementation and not to rank algorithms. Within the reference implementation, doubling the budget to 6,030 evaluations (last column of Table~\ref{tab:spacing-reopt}) raises the mean objective by 5 to 107 units, which is the same order as the recorded LX-SSA--SSA mean differences ($-2$ to $+164$ units). This quantifies why the unequal budgets of Section~\ref{sec:additional-runlevel} prevent any conclusion about relative efficiency and why an equal-budget rerun of the original implementations is required.

